% 「20ワットのコンピュータ」: the popular ("take-along") version, Japanese edition.
% Each chapter paperback/NN.tex corresponds to the full chapter ../chapters/NN_*.tex with the same number;
% on the website the reader switches a chapter between the "Short" and "Full" versions.
\documentclass[11pt,openany]{book}
\usepackage[paperwidth=13cm,paperheight=20cm,top=1.8cm,bottom=1.8cm,inner=1.6cm,outer=1.4cm]{geometry}
% Japanese edition: build with XeLaTeX (see the Makefile)
\usepackage{xeCJK}
\setCJKmainfont[ItalicFont=IPAexGothic,BoldFont=IPAexGothic]{IPAexMincho} % \emph and bold in Gothic: IPAex has no italic
\setCJKsansfont{IPAexGothic}
\setCJKmonofont{IPAexGothic}
\xeCJKsetup{CJKecglue={\hskip 0.15em plus 0.05em minus 0.05em}}
\usepackage{titlesec}
\newcommand{\chaptersuffix}{章}
\titleformat{\chapter}[display]{\normalfont\huge\bfseries}{\ifx\chaptersuffix\empty\appendixname\thechapter\else 第\thechapter\chaptersuffix\fi}{20pt}{\Huge}
\AtBeginDocument{\renewcommand{\contentsname}{目次}\renewcommand{\figurename}{図}\renewcommand{\tablename}{表}\renewcommand{\bibname}{参考文献}\renewcommand{\appendixname}{付録~}}
\usepackage{amsmath,amssymb}
\usepackage{xcolor}
\usepackage[unicode,colorlinks=true,linkcolor=blue!60!black]{hyperref}
\usepackage{microtype}
\input{paperback/pencil} % minimal colour-pencil illustrations, one per chapter
\setlength{\parskip}{0.3em}
\setlength{\emergencystretch}{2em} % narrow paperback page: allow a little extra stretch
\renewcommand{\chaptermark}[1]{\markboth{\small #1}{}}
% neuron type code: first four letters of the primary mediator (as in the full edition)
\newcommand{\nt}[1]{\textsf{\textbf{#1}}}
% pointer to the full edition at the end of each chapter
\newcommand{\fullversion}[1]{\par\bigskip\noindent\small\emph{数式、モデル、問題、出典は完全版の第~#1~章にある。}\normalsize}

\title{\Huge 20ワットのコンピュータ\\[14pt] \Large 脳はどのように計算するか——短く、やさしく}
\hypersetup{pdftitle={20ワットのコンピュータ. 脳はどのように計算するか——短く、やさしく},pdfauthor={Konstantin Kladko}}
\author{\Large コンスタンチン・クラドコ}
\date{}

\begin{document}
\frontmatter
\input{paperback/cover} % cover: a collage of the chapter drawings
% copyright page (same licence as the full edition)
\thispagestyle{empty}
\vspace*{\fill}
\noindent{\footnotesize \copyright~2026 Konstantin Kladko（コンスタンチン・クラドコ）。\\[4pt]
本書の本文と図は、クリエイティブ・コモンズ 表示—非営利—継承 4.0 国際ライセンス（CC BY-NC-SA 4.0）の下で公開されています：\url{https://creativecommons.org/licenses/by-nc-sa/4.0/deed.ja}。著者を表示し、派生物を同じライセンスで公開する限り、非営利目的で自由に複製し、授業に使い、翻訳し、改変することができます。\\[4pt]
完全版、ソース、モデル：\\
\url{https://github.com/kladkogex/20-watt-computer}。}
\clearpage
\tableofcontents
\input{paperback/00_preface}
\mainmatter
\input{paperback/01}
\input{paperback/02}
\input{paperback/03}
\input{paperback/04}
\input{paperback/05}
\input{paperback/06}
\input{paperback/07}
\input{paperback/08}
\input{paperback/09}
\input{paperback/10}
\input{paperback/11}
\input{paperback/12}
\input{paperback/13}
\input{paperback/14}
\input{paperback/15}
\input{paperback/16}
\input{paperback/17}
\input{paperback/18}
\input{paperback/19}
\input{paperback/20}
\input{paperback/21}
\input{paperback/22}
\end{document}
